\chapter{Crude Oil}\label{ch:m3:crude-oil}

At 14:30 New York time on Monday 20 April 2020 the May contract on West Texas
Intermediate crude settled at minus 37.63 dollars a barrel. The session had
opened on Sunday evening at plus 17.73. The contract had never traded below zero
since it was listed in 1983, and it expired the next day at plus 10.01. For about
twenty minutes that afternoon a holder of the contract had to pay tens of dollars
a barrel to be rid of oil it would otherwise have had to collect, in a town in
Oklahoma whose tanks were three-quarters full and whose remaining space was already
spoken for. This chapter explains what a crude benchmark is, how the three that
price most of the world's oil are built, and why the storage behind a benchmark
matters as much as the oil in it.

\section{Crude quality and the benchmarks}

Crude oil is not one commodity but hundreds of grades, each a mixture of
hydrocarbons whose value to a refinery depends on what it can make from it
(\cref{ch:m3:refined-products-and-cracks}). Two numbers summarise most of that
value.

\begin{definition}[API gravity]\label{def:m3:crude-oil:api}
The \emph{API gravity}\index{API gravity} of a crude oil is
\[
\text{°API} = \frac{141.5}{\mathrm{SG}} - 131.5 ,
\]
where $\mathrm{SG}$ is its specific gravity at \qty{60}{\degree F} relative to water.
Water has 10 degrees API; the lighter (less dense) the oil, the higher its API
gravity.
\end{definition}

\begin{definition}[Sweet crude, sour crude]\label{def:m3:crude-oil:sweet}
A crude oil is \emph{sweet crude}\index{sweet crude} if its sulphur content is
low, conventionally below 0.5\% by weight, and \emph{sour crude}\index{sour crude}
if it is higher. (The US Energy Information Administration's statistics use a 1\%
threshold.)
\end{definition}

Light, \omterm{def:m3:crude-oil:sweet}{sweet crudes} yield more gasoline and diesel with less processing and are
worth more; heavy, sour ones need more complex refineries. A crude of specific
gravity 0.827 has 39.6 degrees API. Physical contracts adjust for quality with
escalators: so much per degree API above a reference, so much off per tenth of a
percent of sulphur. With \qty{0.04}{\$/\barrel} per degree and
\qty{0.12}{\$/\barrel} per 0.1\% of sulphur, a 42° crude with 0.25\% sulphur is worth
\qty{0.34}{\$/\barrel} more than a 38° reference with 0.40\%.

\begin{definition}[Benchmark crude]\label{def:m3:crude-oil:benchmark}
A \emph{benchmark crude}\index{benchmark crude} is a crude oil (or a basket of
crudes) whose price is assessed or traded with enough volume and transparency that
the prices of other grades are set as differentials to it.
\end{definition}

Three families of benchmarks price most of the world's oil: the North Sea complex
built around \omterm{def:m3:crude-oil:brent}{Dated Brent}, the US inland barrel WTI delivered at Cushing,
Oklahoma, and the Middle-East benchmarks Dubai and Oman, now joined by Murban.

\section{The North Sea complex}

The seaborne benchmark is not one price but a set of linked markets, from
physical cargoes loading in the next weeks to futures years ahead
(\cref{fig:m3:crude-oil:complex}).

\begin{definition}[Dated Brent, forward Brent, Brent CFD]\label{def:m3:crude-oil:brent}
\emph{Dated Brent}\index{Dated Brent} is the \omterm{def:m3:physical-commodity-markets:pra}{price-reporting agency}'s daily
assessment of physical North Sea cargoes with loading dates already fixed in the
coming weeks, set by the most competitive of a basket of grades after quality
adjustments. \emph{Forward Brent}\index{forward Brent} is the market in forward
contracts for a cargo of one of the basket grades loading in a named future month,
whose grade and loading dates the seller nominates later. A \emph{Brent
CFD}\index{Brent CFD} is a short-dated swap that exchanges the average of Dated Brent
over one week for a fixed price agreed against forward Brent.
\end{definition}

\begin{omfigure}
\centering
\begin{tikzpicture}[font=\small,
  lay/.style={draw,thick,rounded corners=2pt,align=center,minimum height=0.9cm,text width=3.7cm},
  arr/.style={{Stealth[length=2.2mm]}-{Stealth[length=2.2mm]},thick}]
\node[lay,draw=omDef,fill=omDef!8] (dated) at (0,0) {Dated Brent\\physical cargoes,\\dates fixed};
\node[lay,draw=omProp,fill=omProp!8] (cfd) at (4.9,0) {Brent CFDs\\weekly swaps};
\node[lay,draw=omThm,fill=omThm!8] (fwd) at (9.8,0) {forward Brent\\cargo in a named month};
\node[lay,draw=omExo,fill=omExo!8] (fut) at (9.8,-2.2) {ICE Brent futures\\cash-settled on an index};
\draw[arr] (dated) -- (cfd);
\draw[arr] (cfd) -- (fwd);
\draw[arr] (fwd) -- node[right,font=\footnotesize,align=left] {exchange for\\physical} (fut);
\node[font=\footnotesize,align=center] at (3.2,-2.2) {time to loading grows from left to right:\\days and weeks \quad$\longrightarrow$\quad months and years};
\end{tikzpicture}

\omcaption{The Brent complex. The CFDs price the difference between the physical
cargo of the coming weeks and the forward month; an exchange for physical turns a
futures position into a forward cargo. Schematic.}\label{fig:m3:crude-oil:complex}
\end{omfigure}

The basket exists because the output of the original Brent field fell long ago;
more grades make the benchmark harder to squeeze. Each grade enters with a
quality premium or de-escalator, and on each day the lowest net price sets the
assessment. If on one day Forties is offered at \qty{81.05}{\$/\barrel} and the
other grades, net of their premiums, at \qty{81.20}{\$/\barrel} or more, \omterm{def:m3:crude-oil:brent}{Dated Brent}
is \qty{81.05}{\$/\barrel}.

\begin{dated}{2026-09}{The basket}\label{dat:m3:crude-oil:basket}
Platts' \omterm{def:m3:crude-oil:brent}{Dated Brent} reflects Brent, Forties, Oseberg, Ekofisk and Troll and, since June
2023 deliveries, WTI Midland, a US grade delivered into Rotterdam, with assessments
kept on an FOB basis, for cargoes of 700\,000 barrels (up from 600\,000)
loading from ten days to a month ahead. Platts describes the benchmark as the reference
point for around two-thirds of the world's oil trade. ICE Brent futures are cash-settled against the ICE
Brent Index, with an exchange for physical into forward cargoes.
\end{dated}

\section{The US inland benchmark}

WTI is a pipeline crude, and its futures contract is physically delivered at a
single inland hub: Cushing, Oklahoma, a town of tank farms where pipelines from the
producing basins meet the lines to refineries. Because the futures are delivered
there, the storage at Cushing is part of the contract.

\begin{definition}[Last trading day]\label{def:m3:crude-oil:ltd}
The \emph{last trading day}\index{last trading day} of a futures contract is the
last day on which it can be traded. A position still open at its close must be
settled by delivery (or, for a cash-settled contract, at the final settlement
price).
\end{definition}

\begin{dated}{2026-09}{The WTI futures}\label{dat:m3:crude-oil:cl}
NYMEX Light Sweet Crude Oil futures: 1\,000 barrels, delivery \omterm{def:m3:physical-commodity-markets:incoterms}{free on board} at any
pipeline or storage facility in Cushing, Oklahoma. Trading terminates three business
days before the 25th calendar day of the month before the delivery month (three
business days before the last business day preceding the 25th if it is not a business
day). The active month rolls to the next contract some days before expiry, and the
other months then settle from calendar-spread trades.
\end{dated}

Most US pipeline crude does not price on a single day either.

\begin{definition}[Calendar month average]\label{def:m3:crude-oil:cma}
A \emph{calendar month average}\index{calendar month average} (CMA) price is the
average of a benchmark's daily settlements or assessments over the business days of
a calendar month.
\end{definition}

A barrel delivered in a month and priced on the CMA of the futures is exposed, on
each day of that month, to the front contract, which changes identity when it
expires in the middle of the month: CMA formulas are therefore often quoted with a
``roll'' adjustment for the part of the month priced on the second contract.

\section{Middle-East pricing}

The large Gulf producers sell most of their crude on term contracts to refiners,
not in a spot market, at prices they set themselves.

\begin{definition}[Official selling price]\label{def:m3:crude-oil:osp}
An \emph{official selling price}\index{official selling price} (OSP) is the monthly
differential a national oil company publishes for each grade and destination
region, to be added to a named benchmark average over the loading month to price
its term cargoes.
\end{definition}

For Asia the reference has long been the average of Dubai and Oman assessments;
for August 2025 Saudi Aramco set its Arab Light for Asia at \qty{2.20}{\$/\barrel}
above that average. In March 2021 a futures exchange in Abu Dhabi began trading
the first futures on Murban, a light (40° API) grade physically delivered at
Fujairah, and in June 2021 the Abu Dhabi producer moved its Murban OSP to that
contract.

\begin{dated}{2026-09}{The Gulf in 2026}\label{dat:m3:crude-oil:hormuz}
After US and Israeli air strikes on Iran on 28 February 2026, flows of crude and
products through the Strait of Hormuz fell, in the International Energy Agency's
words, ``from around 20~mb/d before the war to a trickle''. The IEA called it the
largest supply disruption in the history of the global oil market. Brent futures
traded ``within a whisker of \$120/bbl'' in March; IEA members agreed on 11 March to
make 400 million barrels of emergency stocks available. The strait was
``effectively closed again in early July''; \omterm{def:m3:crude-oil:brent}{Dated Brent} ended July at
\qty{96.80}{\$/\barrel}. The seaborne benchmark rose far above the inland one
(\cref{fig:m3:physical-commodity-markets:differential}).
\end{dated}

\section{20 April 2020 step by step}

In the first quarter of 2020 an already oversupplied market lost a large share of
its demand to the pandemic. US crude stocks rose by 15.8 million barrels a week for
four weeks. Cushing, with 75.8 million barrels of working storage capacity, went
from about half full in mid-March to about 76\% on 17 April, and press reports said
most of the remaining space was committed. A holder of the May contract at expiry
had to take 1\,000 barrels a contract at Cushing, which needed tank space or pipeline
capacity booked in advance; most holders who were not in that business had to sell
before the close on 21 April.

\begin{proposition}[The storage floor]\label{prop:m3:crude-oil:floor}
If a trader can take delivery of the near contract, store the oil for $m$ months at
cost $c$ per barrel-month and deliver it into the far contract, then, ignoring
financing, the near price $F_1$ and the far price $F_2$ satisfy
$F_1 \ge F_2 - c\,m$. The spread $F_2 - F_1$ can exceed the cost of storage only
if storage cannot be had at that cost.
\end{proposition}

\begin{proof}
Otherwise buy the near contract, take delivery, store, and sell the far contract:
the cash flows are $-F_1 - cm + F_2 > 0$ with no risk.
\end{proof}

On 20 April the June contract settled at \qty{20.43}{\$/\barrel}: the May--June
spread went from $-\$6.76$ on the Friday to $-\$58.06$, the widest in the EIA's
daily history since 1985 (the next widest is \$8.49). A trader who could have stored a
barrel at Cushing for a month would have been paid \$58.06 for doing so, against storage that CME Group
reports had cost about 50 cents a barrel in a Cushing storage auction in March. Nobody with
access to empty tanks was left to take the trade, and the floor of
\cref{prop:m3:crude-oil:floor} did not bind.

\begin{omfigure}
\begin{tikzpicture}
\begin{axis}[width=11cm,height=5cm,
  xlabel={calendar days from 2 March 2020},ylabel={\$/bbl},
  tick label style={font=\small},label style={font=\small},
  xmin=0,xmax=90,ymin=-45,ymax=55,grid=major,grid style={gray!25},
  legend columns=2,legend style={font=\footnotesize,at={(0.5,-0.3)},anchor=north,draw=none,fill=none,/tikz/every even column/.append style={column sep=8pt}},
  table/col sep=comma]
\addplot[omDef,very thick,mark=*,mark size=1pt] table[x=day,y=c1]{figdata/markets-3/02-crude-oil/spring2020.csv};
\addplot[omThm,very thick,dashed] table[x=day,y=c2]{figdata/markets-3/02-crude-oil/spring2020.csv};
\draw[gray] (axis cs:0,0) -- (axis cs:90,0);
\node[font=\footnotesize,anchor=west] at (axis cs:51,-37.63) {$-37.63$ on 20 April};
\legend{front contract (May until 21 April),second contract}
\end{axis}
\end{tikzpicture}

\omcaption{WTI futures, daily settlements, March to May 2020. The front contract
fell below the second as storage filled; on its penultimate day it settled at
$-\$37.63$ while the June contract settled at \$20.43. Data: US Energy Information
Administration, NYMEX contracts 1 and 2.}\label{fig:m3:crude-oil:spring}
\end{omfigure}

\begin{omfigure}
\centering
\begin{tikzpicture}[font=\footnotesize]
\draw[thick,-{Stealth[length=2.5mm]}] (0,0) -- (13,0);
\foreach \x/\t/\l in {0.3/{Sun 18:00}/{session opens, \$17.73},3.0/{Mon 13:00}/{steep fall begins},5.2/{Mon 14:08}/{\$0},7.4/{Mon 14:29}/{low, $-\$40.32$},9.4/{Mon 14:30}/{settles $-\$37.63$},12.0/{Tue 14:30}/{expires, \$10.01}}{
  \draw[thick] (\x,0.12) -- (\x,-0.12);
  \node[anchor=south,align=center] at (\x,0.2) {\t};
  \node[anchor=north,align=center,text width=2cm] at (\x,-0.2) {\l};}
\end{tikzpicture}

\omcaption{The last two sessions of the May 2020 WTI contract, New York time. The
settlement at 14:30 is the average price of calendar-spread trades in the two
minutes before, applied to the June contract. The price fell from \$0 to $-\$39.55$
between 14:08 and 14:28. Source: CFTC interim staff report,
November 2020.}\label{fig:m3:crude-oil:timeline}
\end{omfigure}

Three features of the contract turned a glut into a negative price. It was
physically delivered at one place whose storage was nearly full. Open interest in
the May contract, 634\,727 contracts at the start of April, was still 108\,593
contracts at the start of the 20 April session, higher than usual for the
penultimate day. And the settlement, taken from spread trades in a two-minute
window, came at the end of the session in which those positions were being closed.
Funds that held front-month futures moved away from it: United States Oil Fund,
an exchange-traded oil fund that normally held the front contract, spread its
holdings over later months in the days around the event. The exchange had warned
its members earlier in April that prices could go negative, and it switched its
option models to the Bachelier model (One Quant Book 2, chapter 13), which allows
negative underlying prices. \Cref{fig:m3:crude-oil:spread} shows how
exceptional the day was: in 39 years of daily data the spread exceeded \$5 on only
25 days.

\begin{omfigure}
\begin{tikzpicture}
\begin{axis}[width=11cm,height=4.8cm,
  ylabel={second minus front, \$/bbl},
  tick label style={font=\small},label style={font=\small},
  xmin=1985,xmax=2024.5,ymin=-8,ymax=11,grid=major,grid style={gray!25},
  xtick={1985,1990,1995,2000,2005,2010,2015,2020},xticklabel style={/pgf/number format/1000 sep={}},
  table/col sep=comma]
\addplot[omDef,thin] table[x=t,y=spread]{figdata/markets-3/02-crude-oil/spread.csv};
\draw[gray] (axis cs:1985,0) -- (axis cs:2024.5,0);
\node[font=\footnotesize,anchor=west] at (axis cs:1986,8.3) {20 April 2020: \$58.06 (off scale)};
\end{axis}
\end{tikzpicture}

\omcaption{WTI second contract minus front contract, 1985 to April 2024: Wednesdays,
plus every day above \$10, capped at \$10. Positive values are contango. Data: US
Energy Information Administration, NYMEX contracts 1 and 2.}\label{fig:m3:crude-oil:spread}
\end{omfigure}

\section{Tutorial: the expiry calendar and the storage floor}\label{tut:m3:crude-oil:floor}

\begin{tutorial}
\textbf{Goal.} Compute the \omterm{def:m3:crude-oil:ltd}{last trading day} of any WTI contract and the storage cost
the market implied on 20 April 2020. \textbf{End state:}
\cref{fig:m3:crude-oil:spring,fig:m3:crude-oil:spread}, expiry on 21 April 2020 and
an implied storage cost of \$58.06 per barrel-month.
\begin{tutsteps}
\item \textbf{The calendar.} Business days, then the expiry rule of the dated box.
\omcode{code/firm/crude/firm_crude.py}{33}{54}{Business days and the WTI last trading day.}\label{lst:m3:crude-oil:ltd}
\item \textbf{The floor.} \Cref{prop:m3:crude-oil:floor} and its inverse.
\omcode{code/firm/crude/firm_crude.py}{65}{74}{The storage floor and the implied storage cost.}\label{lst:m3:crude-oil:floor}
\item \textbf{Run} \texttt{m3\_crude.april\_20()}, \texttt{spread\_stats()} and
\texttt{fig\_crude.py} on the EIA series of contracts 1 to 4.
\end{tutsteps}
\textbf{What to change next.} Add the exchange holiday calendar and list every
expiry since 2015; compute the implied storage cost from contracts 2 and 3 in the same
weeks and see how much less extreme it was.
\end{tutorial}

\section{Build: the crude toolkit}\label{bld:m3:crude-oil:crude}

\begin{build}
\bfield{Purpose} The miniature firm prices physical crude against three benchmark
families: it needs quality arithmetic, the most-competitive-grade rule, the expiry
calendar of the futures it hedges with, calendar-month averages and the storage floor
that tells it when a spread is paying for tanks.
\bfield{Interface} \texttt{api\_gravity(sg)}, \texttt{specific\_gravity(api)};
\texttt{quality\_adjusted(price, api, sulphur, ref\_api, ref\_sulphur, per\_api,
per\_tenth\_sulphur)}; \texttt{basket\_benchmark(grade\_prices, premiums)};
\texttt{cl\_last\_trading\_day(year, month, holidays)};
\texttt{calendar\_month\_average}; \texttt{storage\_floor};
\texttt{implied\_storage\_cost}.
\bfield{Rules} Business days are weekdays not in the holiday set; a month with no
settlement is an error; ties in the basket go to the alphabetically first grade, so
the result is deterministic.
\bfield{Acceptance tests} \texttt{code/firm/crude/tests/}: water at 10° API; the
May 2020 contract expiring on 21 April 2020; the January 2024 contract with and without
the Christmas holiday; the 20 April 2020 implied storage cost.
\bfield{Stretch} Feed \texttt{firm.contracts} (One Quant Book 1, chapter 18) with the
expiry dates; a CMA with the roll adjustment for the part of the month priced on the
second contract.
\end{build}

\begin{omsources}
\item CFTC, \emph{Interim Staff Report on Trading in NYMEX WTI Crude Oil Futures
Contract Leading up to, on, and around April 20, 2020}, November 2020.
\item NYMEX Rulebook, chapter 200 (Light Sweet Crude Oil Futures), via a clearing
member's specification page; ICE Brent Crude futures specification.
\item S\&P Global Platts, release of 8 June 2022 on WTI Midland in Dated Brent; \emph{Specifications Guide: Europe and Africa Crude
Oil}, November 2025, and Brent FAQ (Internet Archive copies).
\item CME Group, ``The Peaking of Storage'' (education article, 2020; Internet Archive copy).
\item US Energy Information Administration: glossary (API gravity); \emph{Today in
Energy} on Murban futures (2021); daily NYMEX futures prices, contracts 1--4.
\item IEA, \emph{Oil Market Report}, March and August 2026.
\item United States Oil Fund, Form 8-K filings, April 2020.
\end{omsources}

\section{Exercises}

\begin{exercise}[$\star$]\label{exo:m3:crude-oil:1}
A crude has specific gravity 0.87. What is its \omterm{def:m3:crude-oil:api}{API gravity}? Is it lighter or heavier
than a 39.6° crude?
\end{exercise}

\begin{exercise}[$\star$]\label{exo:m3:crude-oil:2}
Give the \omterm{def:m3:crude-oil:ltd}{last trading day} of the WTI contract for December 2026 (25 November 2026 is
a Wednesday; ignore holidays other than weekends).
\end{exercise}

\begin{exercise}[$\star$]\label{exo:m3:crude-oil:3}
Why is \omterm{def:m3:crude-oil:brent}{Dated Brent} assessed from a basket of grades rather than from Brent alone?
\end{exercise}

\begin{exercise}[$\star\star$]\label{exo:m3:crude-oil:4}
Using the escalators of the text, what is a 36° crude with 0.60\% sulphur worth
against a reference of \qty{80}{\$/\barrel} at 38° and 0.40\%?
\end{exercise}

\begin{exercise}[$\star\star$]\label{exo:m3:crude-oil:5}
Storage at Cushing costs \qty{0.40}{\$/\barrel} a month and money is free. The second
contract is at \qty{25}{\$/\barrel}. What is the floor on the front contract? What does
a front price of \$18 tell you?
\end{exercise}

\begin{exercise}[$\star\star$]\label{exo:m3:crude-oil:6}
Arab Light for Asia is priced at the loading month's Oman/Dubai average plus the OSP.
With an OSP of \qty{2.20}{\$/\barrel} and an average of \qty{71.35}{\$/\barrel}, what does
the buyer pay for a 2-million-barrel cargo?
\end{exercise}

\begin{exercise}[$\star\star\star$]\label{exo:m3:crude-oil:7}
\emph{Coding.} With \texttt{load\_wti}, count the days on which the second contract
exceeded the front by more than \$5, and give the widest and the second-widest spread
in the history.
\end{exercise}

\begin{exercise}[$\star\star\star$]\label{exo:m3:crude-oil:8}
\emph{Find the flaw.} ``Negative prices are impossible for a commodity with any use:
someone will always take free oil.'' Correct it.
\end{exercise}

\section{Problem: The Last Day of the May Contract}

\begin{problem}[{Weekend problem --- 20 April 2020 from the long side}]\label{pb:m3:crude-oil:1}
On the morning of Monday 20 April 2020 a fund holds 1\,000 May WTI contracts. It has
no storage at Cushing and no pipeline space. On the Friday the May contract settled
at \qty{18.27}{\$/\barrel} and the June at \qty{25.03}{\$/\barrel}; normal Cushing storage
leases cost \qty{0.40}{\$/\barrel} a month (illustrative).

\medskip\noindent\textbf{Part I --- The position.}
\begin{enumerate}
\item How many barrels does the fund have to take if it holds to expiry?
\item When is the \omterm{def:m3:crude-oil:ltd}{last trading day}, and why is 20 April the penultimate day?
\item What is the fund's mark-to-market value at Friday's settlement?
\item What was the May--June spread on Friday, and what storage cost did it imply?
\item Was that consistent with the normal lease?
\end{enumerate}

\medskip\noindent\textbf{Part II --- The day.}
\begin{enumerate}[resume]
\item The May contract settles at $-\qty{37.63}{\$/\barrel}$. What is the fund's
one-day variation margin?
\item What storage cost does the May--June spread now imply?
\item How many times the normal lease is that?
\item Why did the storage floor not bind?
\item What would a trader with 1 million barrels of empty tank at Cushing have earned
by buying May and selling June at settlement, storing for a month at the normal cost?
\end{enumerate}

\medskip\noindent\textbf{Part III --- The alternatives.}
\begin{enumerate}[resume]
\item What would the fund have paid to roll into June at Friday's settlements?
\item What does a long-only fund lose from rolling in contango, per barrel, at Friday's spread?
\item Why did holding later months protect a fund?
\item Why did the settlement method (spread trades in a two-minute window) matter?
\item Which model did the exchange switch its options to, and why?
\end{enumerate}

\medskip\noindent\textbf{Part IV --- Judgement.}
\begin{enumerate}[resume]
\item Who was on the other side: who can buy a contract at $-\$37.63$?
\item What does the day say about the difference between a cash-settled and a
physically settled benchmark?
\item Would a Brent futures contract have gone negative on the same day? Why?
\item State the \emph{named result}: the storage cost the market implied on 20 April.
\item In one sentence: what is a physically delivered futures price a price of?
\end{enumerate}
\end{problem}

\section{Interview questions}

\begin{interviewq}[$\star$ \iqroles{trader}]\label{iq:m3:crude-oil:1}
What is the difference between \omterm{def:m3:crude-oil:brent}{Dated Brent} and the Brent futures price?
\end{interviewq}

\begin{interviewq}[$\star$ \iqroles{trader, researcher}]\label{iq:m3:crude-oil:2}
How can the price of a commodity be negative?
\end{interviewq}

\begin{interviewq}[$\star\star$ \iqroles{researcher}]\label{iq:m3:crude-oil:3}
You backtest a front-month WTI strategy on a continuous series and find a trade on 20
April 2020 that made a fortune. What do you check?
\end{interviewq}

\begin{interviewq}[$\star\star$ \iqroles{trader}]\label{iq:m3:crude-oil:4}
The front WTI spread is $-\$1.50$ and a month of Cushing storage costs \$0.40. What
trade does this suggest, and what can go wrong?
\end{interviewq}

\begin{interviewq}[$\star\star$ \iqroles{developer, risk}]\label{iq:m3:crude-oil:5}
Your pricing library computes option prices with the Black formula on the log of
the futures price. What breaks when the futures go negative, and what do you do?
\end{interviewq}

\begin{interviewq}[$\star\star\star$ \iqroles{trader, risk}]\label{iq:m3:crude-oil:6}
Why does a benchmark based on a basket of grades and several markets resist squeezes
better than a single-grade physical contract, and what does it cost?
\end{interviewq}
